%% kapa: exam=1 round=33 year=2022 period=2 form=A subject=appraisal_law subject_ko=감정평가관계법규 seq=16 number=16
%% source: Q-Net 감정평가사 기출문제 (cst00309) artlSeq=5209997 file=raw/1008/5209997/2022년도 제33회 감정평가사 1차시험 2교시(감정평가 관계법규, 회계학).pdf
%% conversion: kordoc 4.24.51 -> tools/gen1.py (auto); review=proofread-sonnet; flags=none
\begin{question}{16}
감정평가 및 감정평가사에 관한 법령상 감정평가법인등에 관한 설명으로 옳지 않은 것은?
\begin{choices}
  \item 부정한 방법으로 감정평가사의 자격을 받았다는 사유로 감정평가사 자격이 취소된 후 1년이 경과되지 아니한 사람은 감정평가법인등의 사무직원이 될 수 없다.
  \item 감정평가법인은 국토교통부장관의 허가를 받아 토지등의 매매업을 직접 할 수 있다.
  \item 감정평가법인등이나 그 사무직원은 업무수행에 따른 수수료와 실비 외에는 어떠한 명목으로도 그 업무와 관련된 대가를 받아서는 아니 된다.
  \item 감정평가사가 고의 또는 중대한 과실 없이 감정평가서의 적정성을 잘못 심사한 것은 징계사유가 아니다.
  \item 한국감정평가사협회는 감정평가를 의뢰하려는 자가 해당 감정평가사에 대한 징계 사실을 확인하기 위하여 징계 정보의 열람을 신청하는 경우에는 그 정보를 제공하여야 한다.
\end{choices}
\end{question}
\answer{2}
